% =====================================================================
% Chapter 7 -- Uncertainty, Calibration, and Explainability
% =====================================================================
\documentclass[../preamble.tex]{subfiles}
\begin{document}

\chapter{Uncertainty, Calibration, and Explainability}
\label{ch:trust}

\section{Operating-point selection}
\label{sec:trust-op}

The default decision threshold of 0.5 is rarely the right choice for an
imbalanced, screening-oriented classifier. The thesis therefore computes
three operating points:

\begin{itemize}
  \item the default $0.5$ threshold,
  \item the Youden-J operating point, which maximises $\mathrm{sensitivity} +
        \mathrm{specificity} - 1$, and
  \item the sensitivity-targeted operating point, which fixes
        $\mathrm{sensitivity} = 0.95$ and reads the resulting threshold.
\end{itemize}

The threshold-sweep figure in \cref{fig:threshold-sweep} shows the full
sensitivity / specificity / F1 curves, with the three operating points
highlighted.

\begin{figure}[htbp]
  \centering
  \includegraphics[width=\linewidth]{figures/thesis_threshold_sweep.png}
  \caption{Threshold sweep for the pre-HPO ensemble on the held-out PCOSgen
           test set. The Youden-J threshold is 0.6198 and the
           sensitivity-targeted threshold is 0.7597.}
  \label{fig:threshold-sweep}
\end{figure}

\section{Calibration}
\label{sec:trust-calib}

Neural networks are systematically overconfident on classification tasks
\cite{guo2017}. Temperature scaling, which divides the logits by a learned
scalar $T$, preserves the original accuracy while improving the calibration
error. The thesis applies temperature scaling in two passes:

\begin{enumerate}
  \item \textbf{Pass 1.} For each model, learn $T_i$ on the validation set
        using
        \begin{equation}
          T_i = \arg\min_T \mathrm{NLL}(\sigma(z_i / T),\, y_i),
          \label{eq:temp-each}
        \end{equation}
        where $z_i$ are the logits, $y_i$ are the labels, and $\sigma$ is the
        softmax operator.
  \item \textbf{Pass 2.} Average the calibrated probabilities across the
        three ensemble members, then learn a single $T_{\mathrm{ens}}$ on the
        same validation set.
\end{enumerate}

The two-pass procedure is summarised in \cref{alg:two-pass-calib}. The
per-model and ensemble temperatures, the intermediate ECE, and the final ECE
are recorded in \cref{tab:calibration-detail}.

\begin{algorithm}[htbp]
  \caption{Two-pass temperature scaling for the ensemble.}
  \label{alg:two-pass-calib}
  \begin{algorithmic}[1]
    \State Inputs: validation logits $\{z_i\}$, validation labels $\{y_i\}$,
                   ensemble members $\{m_1, m_2, m_3\}$.
    \For{each $m_k \in \{m_1, m_2, m_3\}$}
      \State $T_k \gets \arg\min_{T \in [0.1, 10]} \mathrm{NLL}(\sigma(z_k / T), y_k)$.
    \EndFor
    \State $p_{\mathrm{avg}} \gets \frac{1}{3}\sum_{k=1}^{3} \sigma(z_k / T_k)$.
    \State $T_{\mathrm{ens}} \gets \arg\min_{T \in [0.1, 10]} \mathrm{NLL}(p_{\mathrm{avg}}^{1/T}, y)$.
    \State Return calibrated probabilities
            $\sigma(z_k / T_k)$ per member and $p_{\mathrm{avg}}^{1/T_{\mathrm{ens}}}$
            for the ensemble.
  \end{algorithmic}
\end{algorithm}

\begin{figure}[htbp]
  \centering
  \includegraphics[width=\linewidth]{figures/paper_fig_calibration_compare.png}
  \caption{Per-model calibration before and after temperature scaling. See
           \cref{tab:calibration-detail} for the exact values.}
  \label{fig:calibration}
\end{figure}

\begin{table}[htbp]
  \centering
  \small
  \caption{Two-pass temperature scaling. Per-model temperatures scale the logits of each member before probability averaging; the ensemble temperature then scales the averaged logits.}
  \label{tab:calibration-detail}
  \begin{tabular}{lr}
    \toprule
    \textbf{Component} & \textbf{Value} \\
    \midrule
    densenet121 \,\,$T_i$ & 2.0162 \\
    convnext_tiny \,\,$T_i$ & 1.9579 \\
    vit_base \,\,$T_i$ & 2.1982 \\
    Ensemble $T_\text{ens}$ & 0.9306 \\
    ECE after pass 1 & 0.0854 \\
    ECE after pass 2 & 0.0810 \\
    Test samples & 734 \\
    \bottomrule
  \end{tabular}
\end{table}


The reduction in ECE is modest in absolute terms (0.0854 to 0.0810) but the
calibration plot shows a meaningful alignment of predicted probabilities
with observed accuracy. All three fine-tuned members and the ensemble are
overconfident before temperature scaling; the ensemble temperature is the
smallest of the four values, which is consistent with averaging reducing
individual-model variance.

\section{Uncertainty}
\label{sec:trust-unc}

Predictive uncertainty is estimated using MC-Dropout with dropout active at
inference time \cite{gal2016}. For each test image, 50 stochastic forward
passes are performed and the predictive entropy is computed:
\begin{equation}
  \mathbb{H}(\hat{y}) = -\sum_{c \in \{0,1\}} \bar{p}_c \log \bar{p}_c,
  \label{eq:entropy}
\end{equation}
where $\bar{p}_c$ is the mean posterior probability of class $c$ across the
50 passes.

\begin{figure}[htbp]
  \centering
  \includegraphics[width=\linewidth]{figures/thesis_uncertainty_hist.png}
  \caption{MC-Dropout entropy distributions for the three ensemble members,
           showing a clear separation between correctly and incorrectly
           classified cases. The mean entropies are 0.234 nats (DenseNet-121),
           0.215 nats (ConvNeXt-Tiny), and 0.160 nats (ViT-B/16).}
  \label{fig:uncertainty-hist}
\end{figure}

The ViT-B/16 model has the lowest mean entropy (0.160 nats), suggesting it
remembers more confident posteriors; the DenseNet-121 has the highest mean
entropy (0.234 nats). The median entropy follows the same ordering. The
broader distribution of correct-case entropies for the convolutional models
indicates that uncertainty is not a simple proxy for correctness, but a
useful complementary signal.

\section{Uncertainty-gated referral as a design concept}
\label{sec:trust-referral}

The original FYDP brief proposed an uncertainty-gated clinical referral
mechanism: route the patient to a clinician not only when the model
predicts PCOS but also when the model is uncertain. This thesis implements
that mechanism as a design concept and reports the per-coverage cohort
statistics, but does not deploy the mechanism in a clinical workflow. The
trade-off between auto-decided coverage and accuracy on the auto-decided
subset is recorded in the per-checkpoint \texttt{mc\_dropout\_results.json}
files. A deployed version would require a prospective study, a defined
confidence threshold, and a calibrated human-cost function.

\section{Explainability}
\label{sec:trust-xai}

The thesis uses Grad-CAM to expose the regions of an image that most
influenced the prediction \cite{selvaraju2017}. Grad-CAM is computed on the
final convolutional layer of each backbone, so the spatial resolution is
limited to the receptive field of the deepest feature map. For ViT-B/16,
whose "final convolutional layer" is a transformer block, the hook is
attached to the last normalisation layer.

\begin{figure}[htbp]
  \centering
  \includegraphics[width=\linewidth]{figures/paper_fig_xai_grid.png}
  \caption{Grad-CAM grid for the three ensemble members. Each row shows
           representative correct, uncertain, and incorrect cases.}
  \label{fig:xai-grid}
\end{figure}

The qualitative observation is that on cases the model gets correct, the
Grad-CAM signal is concentrated on the central ovarian region. On cases
the model gets wrong, the activation tends to spill onto the image border or
the bottom-corner region, which is consistent with the letterbox-padding
shortcut described in \cref{ch:preprocessing}. After the padding shortcut is
removed, this border-driven behaviour is greatly reduced.

\section{Summary}
\label{sec:trust-summary}

The calibration, uncertainty, and XAI components are complementary views
of the same trained model. Calibration adjusts the model's confidence,
MC-Dropout separates high-confidence from low-confidence predictions, and
Grad-CAM provides a region-level explanation. Together they address the
trust gap that aggregate metrics alone cannot fill.

\end{document}